%% APPENDIX D -- the false-alarm expectation, the control ensemble, and the
%% execution-block quality criterion.  Carries \label{app:falsealarm}
%% (04_method.tex, 05_results.tex, 08_backmatter.tex) and
%% \label{sec:blockquality} (04_method.tex, 05_results.tex).
%%
%% ★★★★ WHAT THIS APPENDIX IS FOR, AND WHAT IT MAY NOT GROW BACK INTO.  Four
%% things and no more: the expectation's definition, the resampling the count
%% is compared through, the principal numbers, and the control model's
%% limitations.  Everything else that stood here is either stated once in
%% \S\ref{sec:trigunif} -- the extraction geometry, the two probabilities and
%% their correction, the bound's qualifications -- or is a diagnostic that
%% decided nothing: the threshold scan's margins and its counts at the next
%% and median blocks, the unscaled form of the block rule, the cell-grain
%% reading of the same expectation, the clean coarse windows' addition to it,
%% the systemic-velocity sensitivity, the near-band ratio read half-level by
%% half-level and its value at the trigger itself, the reserved blocks that
%% share no session, and the scrambling control.  Each was measured and each
%% is in the generators; none of them changes a disposition, and a limitation
%% is condensed here, never deleted.
\section{The false-alarm expectation and the control ensemble}
\label{app:falsealarm}

\subsection{The expectation, and what it is conditional on}

\emph{One attribution set serves both sides of the comparison}: the mask of
Appendix~\ref{app:linemask}, which holds no sulphur monoxide, so the
\SoNCrossWord{} crossings on SO transitions are unattributed and a control
maximum on one is an unattributed exceedance. Each of the \CeNWin{} Class~A
windows of the primary census contributes the fraction of its own \NCtrl{}
control positions that reaches the trigger -- no cell count enters it and no
form is assumed for the tail -- and summed over the windows \CeExpTrig{} are
expected to trigger by chance. How many of those would be attributed is
measured on the channels the control maxima land on. The channel of a maximum exists only for the
\RcNRetCtrl{} control positions per window that retain a dynamic spectrum
(\S\ref{sec:vistest}), so the share is measured on the \CeNExcPos{} such
positions of \CeNExcWin{} of the \CeNWin{} windows, which carry
\CeExcCovPct{} per cent of the expectation above. The full ensemble holds
\CeNExcFull{} exceedances, which divided by \NCtrl{} is the \CeExpTrig{}
above; the \RcNRetCtrl{} are indistinguishable from a uniform draw from
their window's own \NCtrl{} by rank ($P=\CeRetKsP$) and their exceedance
rate from the full ensemble's by \CeRetZ{} standard deviations.
\CeNExcAttr{} of \CeNExc{} exceedances fall inside the mask, carried to the
star's frame exactly as a stellar cell, so the attribution takes \CeMaskPct{}
per cent of the expectation where a uniform landing would take \CeUnifPct.
The atmospheric screen (Appendix~\ref{app:rfi}) then takes \CeTlCtrlCut{} of
\CeTlCtrlSurv, leaving \CeExp{} expected unattributed Class~A crossings
against \CeObs{} observed, a sum over rows with nothing carried at an assumed
rate. Both sides are read over the whole primary beam \S\ref{sec:statistic}
defines, with each block on its own antenna geometry; the narrower field of
\S\ref{sec:trigunif} disagrees with it on \FrNFlag{} of the \CeNWin{} windows
and in both directions: the whole beam holds a crossing the narrower field
does not toward \FrUpStars, both attributed to carbon monoxide, and the
narrower field one the whole beam does not toward \FrDownStars. System by
system \MtNSysLowP{} of \MtNSysChance{} fall below one chance in twenty
against \MtNSysLowPExp{} expected; the \CeNWinTail{} untabulated archival-tail
windows, with the \CeNTailCrossWord{} unattributed crossings toward
\CeTailStar, enter neither side.

\emph{A count is compared with the distribution of counts chance delivers.}
Over \CeNBoot{} resamplings -- each window's rate and attributed share redrawn
from its own probes, and the execution blocks resampled because windows share
them -- a count of \CeObs{} or more arises with probability \CeProbNull{},
conditional on the control-position model and not on an independently
established null. Added to each resampled rate, \CeAddLim{} further crossings
still leave \CeObs{} credible at \CeProbAtLim{} and one more does not
(\CeProbAboveLim{} against \CeAlphaPct{} per cent). The uncertainty in
\CeExp{} is carried inside that bound by the same resampling.

\emph{The one assumption the searched data cannot test is that a star behaves
like a control position in its own field.} The star carries \CeNear{} times a
control position's cells in the band $\CeNearBandLo\le T<\CeNearBandHi$
(\CeNearIntLo--\CeNearIntHi{} at \CeCovPct{} per cent, blocks resampled).
Calibrating the stellar rate against that ratio instead, the sensitivity
test of \S\ref{sec:trigunif}, lowers the bound to \CeAddLimNear; the
whole-beam conclusion that there is no significant excess is unchanged under
it, and the control positions' own rate is adopted as the smaller of the two
and the more demanding of the count under test. The \CeObs{} would stand
above chance at \CeAlphaPct{} per cent only if stars triggered at \CeRCrit{}
times a control position's rate or less, and every ratio measured is an
excess in the census. The reserved blocks trigger in \StHoObs{} of \StHoNFine{} windows
against \StHoExpFine{} predicted, \HsShareObs{} of them in blocks sharing a session with a
searched one, and the \HsIndBlocks{} sharing none hold \HsIndObsWord{}
trigger \HsIndObsNoun{} against \HsIndExp{} expected, so they neither confirm nor contradict the ratio
(Appendix~\ref{sec:heldout}).

Near a band edge the exceedance rate is $\times\VnEdgeRatio$ the interior
rate, pooled -- and the excess is one object: the window along the sightline
to HD~48370 contributes \VnEdgeDomPct{} per cent of every near-edge
exceedance in the survey, without it the ratio is
$\times\VnEdgeNoDomRatio$, and on a held-out half it is
$\times\VnEdgeHalfCleanRatio$ ($\pv{\VnEdgeHalfCleanP}$). No edge
correction is therefore applied to any limit.

\subsection{Execution-block quality}
\label{sec:blockquality}

The comparison assumes each block's control ensemble behaves like noise, and
one block makes that false: all \NCtrl{} of its probes exceed the trigger, the
smallest at \BqFailCtrlMin. Every block is therefore judged on the median of
its windows' control ensembles divided by its class median, which also flags
\DqNWin{} of the \NWindows{} windows individually at \BqFactor{} the survey
median. \emph{The criterion is post hoc}, written after the condition was
seen; it is trustworthy because the failing block sits at \BqFailQ{} and the
next at \BqNextQ, a factor \QsRatio{} apart, and all \QsNStep{} thresholds
between select the same block. \BqNFail{} of \BqNBlocks{} census blocks fails;
its \BqFailNWin{} coarse-channel windows hold the entire Class~B crossing
population, \EvNCrossFlag{} of the \EvNCrossRaw, listed in the ledger and
entering no count elsewhere.

\subsection{What the ensemble can resolve, and in what unit}

Ranking the star against \NCtrl{} control positions decides which windows
deserve examination, but $1/(\NCtrl+1)=\RankFloorVal$ is the ensemble's
\emph{resolution}, not the probability that noise produced the window: a
survey of \NWindows{} windows would need a rank \BonfRankRatio{} times finer,
and correlated controls inside the beam leave an annulus resolving no better
than \RankEffLo--\RankEffHi. That floor is the design's central limitation:
no window can be significant on the rank alone, which is why confirmation
requires independent recurrence.

The widened recurrence window of \S\ref{sec:etacrv} is counted in windows, a
window covering several crossings being one measurement, so the \WuNRead{}
pairings are \WuNWin{} distinct windows, each one's chance formed from its
independent channel--drift cells (\WuExpWinCtrl{}, $P=\WuProbWinCtrl$, with
its control positions supplying the rate instead; the larger $P$ is quoted).
Of the \WuNWinExc{} windows above the trigger,
\WuNWinExcTwo{} are toward \WuTwoStars{} alone, and
over the other \WuNWinRest{} windows the comparison is \WuNWinRestExc{}
against \WuExpRest{} at $P=\WuProbRest$; those two hold \TwNWinTwo{} of the
windows and \TwExpTwo{} of the expectation, and \S\ref{sec:twostars} splits
the departure between them. The full treatment is in
\texttt{documentation/\allowbreak appendix\_D\_false\_alarm.md} of the data archive
\citep{SIdeposit}.

%% The chance expectation here is summed from the measured control exceedance
%% rates and each window's own masked grid share (chance_v414.py, round 310),
%% so no cell count and no Gaussian tail enters it and nothing in it is
%% conditioned on the crossings; the cell-count route is a cross-check only,
%% and the rank-conditioned expectation (\ClustStageMean against
%% \ClustStageNObs) is a pipeline diagnostic that must not be reachable as
%% "the" expectation.  tab:ctrlsep, tab:excess, tab:arraysplit and
%% fig:ctrldiag are generated and nothing cites them; r15-deposit, 2026-10-08:
%% they used to be described as "in the data release", and there is no data
%% release, so they are simply unpublished diagnostics.
\subsection{The registration of the predicted cell}
\label{sec:regbudget}

The recurrence test of \S\ref{sec:etacrv} reads one channel at one drift,
registered between epochs to \RgBudgetChMed{} of a channel in every window,
the frame correction moving it by up to \ZtFrameChanMax{} channels. Widening
each reading to the worst cell within \RgTolKms\,km\,s$^{-1}$, pinning it at
every anchor along the discovery track, and raising the bar to \ZtTrialThr{}
to hold one false confirmation over all \RcNRepeatMeas{} readings each
changes no disposition. What the
recoveries establish about the registration of a single cell rests on the
\RgRegNUnres{} recovery readings of \RgRegStar{} alone, whose line is confined
to one channel and which register to \RgRegMaxCh{} of a channel,
\RgRegMaxKms\,km\,s$^{-1}$; the rest test the frame correction only to the
width of a resolved line, up to \RgResFwhmChMax{} channels.

\subsection{From the widened window's exceedances to the pairs}
\label{sec:pairacct}

Two counts of \S\ref{sec:etacrv} are counts of different objects. The
\WdNExc{} above-trigger readings are window readings, one per pairing of an
unattributed crossing with a window covering its frequency in another block
-- \WdNCov{} such pairings over the \PaNCovWinCross{} crossings a later
catalogued block covers, row 2.3.3 of Table~\ref{tab:counts} -- so a pair of
unattributed crossings is read twice: \PaNPairTwice{} readings are second
readings of \PaNPairRead{} pairs, \PaNPairBeyond{} outside the interval the
pair test admits and \PaNPairInside{} inside it, the set of
Table~\ref{tab:pairs}, the other \TjNCandLost{} of the \TjNCandPair{} wanting
a discovery spectrum or a fitted drift. The accounting for each pair is in
the data archive \citep{SIdeposit},
\texttt{candidates/\allowbreak drift\_pair\_accounting.csv}. A
catalogued window covers \WdNOrbTot{} of the \RcNExcl{} crossings that carry an
amplitude bound and spans the orbital excursion of \S\ref{sec:etacrv} in
\EcNSpan; for \EcNShort, toward \EcShortStarList, the widest covering window is
narrower, and \WdNNoCov{} -- \WdNoCovStar{} -- is covered only by the
uncatalogued tail block \texttt{\WdNoCovEb}, \WdNoCovSepD\,d later. Of the
\TjNPower{} pairs whose epochs are close enough together for the
predicted offset to be tighter than the searched window, every one fails a
drift changing at a constant rate, by \TjSigMin{} times that resolution in
the mildest case. The control row is a line that does not drift: the
\TjCtlLine{} line of \TjCtlStar{} is recovered at one stellar-frame frequency
\TjCtlDtD{} days apart, within \TjCtlNSig{} of the comparison's own
resolution of the prediction, so between these two blocks it resolves a mean
drift to \TjCtlDriftResHzS\,Hz\,s$^{-1}$: a grid step carried over the gap,
which is why a longer gap does not sharpen it.

$\beta$~Pictoris is the closest case, and it is the one the comparison
cannot dismiss; its orbit is argued here and nowhere else, the methanol
identification of its two crossings being excluded in
Appendix~\ref{app:linemask}. Its two crossings, \TjBpDvKms\,km\,s$^{-1}$ apart
in the star's frame in blocks \TjBpDtH\,h apart, have negative fitted drifts
($\TjBpDriftLo$ and $\TjBpDriftHi$\,Hz\,s$^{-1}$), so a single carrier would
have had to rise in frequency between the epochs and under constant
acceleration misses by \TjBpNSig{} times that resolution -- but the mean
acceleration it needs, \TjBpAccelReq\,m\,s$^{-2}$, is inside the
\WdAGrid\,m\,s$^{-2}$ ceiling. The pair fixes three numbers and a circular
orbit about a \EcMassBp\,M$_\odot$ star \citep{Crifo1997} has three
parameters: one orbit (\ObDayBp\,d, \ObIncBp$^\circ$, \ObAuBp\,au) reproduces
all three, eccentric ones fit to $e\approx\ObEccMaxBp$, and the pair is not
separable.

%% ★★ r27-cut-results (row 34, R2-D8): Table D1 (tab:pairacct) MOVED TO THE
%% DEPOSIT as candidates/drift_pair_accounting.csv, generated from
%% pairacct_r23.json; the one-line pointer is in the paragraph above.  The
%% float's removal renumbers the appendix tables after it; every citation is
%% a \ref, so nothing prints a stale number.  tab_pairacct_r23.tex is still
%% generated and is now uncited: retire it from cleanregen.PRODUCTS or keep it
%% as the deposit emitter's source, but do not \input it.

%% ★★ v4.23: THIS SUBSECTION AND ITS TABLE SIT AT THE END OF THE
%% APPENDIX FOR A REASON THAT IS NOT STYLE.  Placed next to the
%% registration budget, where it belongs by subject, its float took
%% the number Table~D1 and pushed the crossing ledger's sample to D2
%% and the recurrence positive controls to D3 -- and the positive
%% controls are the table Referee 2's first point is about, which he
%% and the response letter both call Table~D2.  Renumbering the
%% table a referee is reading about is a cost with no benefit, so the
%% new float goes last and takes the new number.  \S\ref{sec:edge}
%% is cited from \S4 and \S5.6 and the pointer is unaffected.
\subsection{The drift grid's edges}
\label{sec:edge}

The drift trials run from $-\dot\nu_{\max}$ in half-channel steps to the
first trial at or beyond $+\dot\nu_{\max}$, so the searched interval is
asymmetric by up to one step; \EgNEdge{} of the \EgNDrift{} crossings with a
fitted drift lie on an outermost trial (Table~\ref{tab:edge}), toward
\EgEdgeStars, and are reported as lower bounds, \EgNEdgeOut{} of them carrying
no amplitude bound, which is why the table reaches beyond
\S\ref{sec:etacrv}'s range.

\begin{table}
%% ★ v4.26 (r25-captions): 6 printed lines -> the author's 3.  DELETED after
%% checking: "a maximum at the boundary of the interval searched is a lower
%% bound" is \S\ref{sec:etacrv}'s own definition of \emph{edge-limited}, and
%% "Table~\ref{tab:pairs} gives them in the star's frame" is now a clause of
%% that same paragraph, which is where a reader meets both tables.
\caption{The crossings whose fitted drift lies on an outermost trial of
their window's own grid, with that grid. Rates are in the observatory
frame.}
\label{tab:edge}
\footnotesize
%% GENERATED by edge_v423.py -- do not hand-edit.
\begin{tabular}{@{}llr@{~}r@{~}c@{~}r@{~}l@{}}
\hline
Star & block & $\nu$ & $\dot\nu_{\rm obs}$ & trial & \multicolumn{2}{c}{$|a_{\rm los}|$} \\
 & & (GHz) & (Hz\,s$^{-1}$) & & \multicolumn{2}{c}{(m\,s$^{-2}$)} \\
\hline
BD+05~1668 & \texttt{X28a8} & 242.195 & $+3452$ & 2/2 & $\ge$ & 4.27 \\
HD~14055 & \texttt{X29ae} & 331.156 & $-3976$ & 0/135 & $\ge$ & 3.60 \\
HD~14055 & \texttt{X2c3bb} & 345.526 & $+4184$ & 141/141 & $\ge$ & 3.63 \\
HD~161868 & \texttt{X10ca8} & 345.634 & $+4195$ & 127/127 & $\ge$ & 3.64 \\
TRAPPIST$-$1 & \texttt{X259b} & 346.245 & $-4623$ & 0/137 & $\ge$ & 4.00 \\
\hline
\end{tabular}

\end{table}
